\documentclass[french]{article}

\usepackage[utf8]{inputenc}
\usepackage[T1]{fontenc}
\usepackage{lmodern}
\usepackage{geometry}
\usepackage{graphicx}
\usepackage{babel}
\usepackage{amsmath}
\usepackage{amsthm}
\usepackage{amssymb}
\usepackage{hyperref}
\usepackage{stmaryrd}
\usepackage{enumitem}
\usepackage{tcolorbox}
\usepackage{dsfont}
\usepackage{multirow}
\usepackage{etoc}
\usepackage{titlesec}
\usepackage{esint}
\usepackage{cancel}
\usepackage{mathdots}


\setlist[itemize]{label=$\bullet$}


\renewcommand{\P}{\mathbb{P}}
\newcommand{\N}{\mathbb{N}}
\newcommand{\Z}{\mathbb{Z}}
\newcommand{\Q}{\mathbb{Q}}
\newcommand{\R}{\mathbb{R}}
\newcommand{\C}{\mathbb{C}}
\newcommand{\K}{\mathbb{K}}
\renewcommand{\L}{\mathbb{L}}
\newcommand{\U}{\mathbb{U}}
\newcommand{\st}{^*}
\newcommand{\tq}{\middle|}
\newcommand{\inv}{^{-1}}
\newcommand{\E}{\mathbb{E}}
\newcommand{\F}{\mathbb{F}}
\newcommand{\V}{\mathbb{V}}
\renewcommand{\ker}{\text{Ker}}
\newcommand{\im}{\text{Im}}
\newcommand{\card}{\text{Card}}
\newcommand{\Tr}{\text{Tr}}
\newcommand{\id}{\text{id}}

\theoremstyle{plain}
\newtheorem{thm}{Théorème}
\def\thmautorefname{théorème}
\newtheorem{lem}{Lemme}
\def\lemautorefname{lemme}
\newtheorem{prop}{Proposition}
\def\propautorefname{propriété}
\newtheorem{cor}{Corollaire}
\def\propautorefname{corollaire}

\theoremstyle{definition}
\newtheorem{dfn}{Définition}
\def\dfnautorefname{définition}
\newtheorem{exa}{Exemple}
\newtheorem{rmq}{Remarque}

\title{Décompositions et applications}

\begin{document}

\maketitle
\section{Décomposition(s) de Dunford}
\begin{thm}[Décomposition de Dunford]
	Si le polynôme caractéristique de $u$ est scindé sur $\K$, alors il existe un unique couple d'endomorphismes $(d,n)$ tel que
	\begin{enumerate}
		\item $u=d+n$ ;
		\item $d$ soit diagonalisable ;
		\item $n$ soit nilpotente ;
		\item $u$ et $d$ commutent.
	\end{enumerate}
	De plus, de tels endomorphismes sont des polynômes en $u$.
\end{thm}

\begin{cor}
	Si le polynôme caractéristique de $M$ est scindé sur $\K$, alors $M$ est diagonalisable si, et seulement si, $\exp(M)$ l'est.
\end{cor}

\begin{rmq}
	En pratique, pour obtenir la décomposition de Dunford de $u$, on réalise la décomposition en éléments simples de $\displaystyle\frac{1}{\pi_u}$ puis on calcule les projecteurs spectraux $\pi_k$. On pose alors $d=\displaystyle\sum_{k=1}^{p}\lambda_k\pi_k$ où les $\lambda_k$ sont dans $\text{Sp}_\K(u)$ puis on pose $n=u-d$.
\end{rmq}

\begin{cor}
	Si $\chi_A$ est scindé sur $\K$ et $D+N$ est la décomposition de Dunford de $A$, alors la décomposition de Dunford de $\exp(A)$ est $\exp(D)+\exp(D)(\exp(N)-I_n)$.
\end{cor}

\begin{thm}
	Les assertions suivantes sont équivalentes :
	\begin{enumerate}[label=\roman*)]
		\item $\exp(A)=I_n$ ;
		\item $A$ est diagonalisable et $\text{Sp}_\C(A)\subset 2i\pi\Z$.
	\end{enumerate}
\end{thm}

\begin{rmq}
	La décomposition de Dunford est principalement utilisée pour calculer des puissances ou des exponentielles de matrices.
\end{rmq}

\begin{thm}[Décomposition de Dunford multiplicative]
	Si $A\in\mathcal{GL}_n(\K)$ et $\chi_A$ est scindé sur $\K$, alors il existe un unique couple de matrices $(T,D)\in\mathcal{GL}_n(\K)^2$ tel que :
	\begin{enumerate}
		\item $A=TD=DT$ ;
		\item $T$ soit diagonalisable ;
		\item $D$ soit unipotente.
	\end{enumerate}
\end{thm}

\section{Réduction de Jordan}

\begin{dfn}[Réduite de Jordan]
	On appelle réduite de Jordan toute matrice de la forme : $$J_r(\lambda)=\begin{pmatrix}
		\lambda&1&&&(0)\\
		&\lambda&1&& \\
		&&\ddots&\ddots& \\
		&&&\lambda&1 \\
		(0)&&&&\lambda
	\end{pmatrix}\in\mathcal{M}_r(\K)$$
\end{dfn}

\begin{thm}[Réduction de Jordan]
	Si $\chi_A$ est scindé, alors $\pi_A$ est scindé et $A$ est semblable à une matrice diagonale par blocs dont les blocs sont des réduites de Jordan. De plus, cette décomposition est unique à l'ordre près des blocs.
\end{thm}

\begin{cor}
	Deux matrices de $\mathcal{M}_n(\K)$ dont les polynômes caractéristiques sont scindés sont semblables si, et seulement si, elles admettent la même forme de Jordan.
\end{cor}

\begin{cor}
	Si $\chi_A$ est scindé, alors $A$ et $A^T$ sont semblables.
\end{cor}

\section{Décomposition de Frobenius}

\section{Générateurs de $\mathcal{GL}_n(\K)$ et de $\mathcal{SL}_n(\K)$}

\begin{dfn}[Matrices de transvection]
	On appelle matrice de transvection toute matrice qui s'écrit sous la forme $T_{i,j}(\lambda)=I_n+\lambda E_{i,j}$ avec $i\ne j$ et $\lambda\in\K$.
\end{dfn}

\begin{dfn}[Matrices de dilatation]
	On appelle matrice de dilatation toute matrice qui s'écrit sous la forme $D_{i}(\lambda)=I_n+(\lambda-1)E_{i,i}$ avec $\lambda\in\K\st$.
\end{dfn}

\begin{thm}
	Les matrices de transvection et de dilatation engendrent $\mathcal{GL}_n(\K)$. On peut n'utiliser qu'une seule matrice de dilatation.
\end{thm}

\begin{cor}
	Les matrices de transvection engendrent $\mathcal{SL}_n(\K)$.
\end{cor}

\begin{cor}
	Les groupes $\mathcal{SL}_n(\R)$, $\mathcal{SL}_n(\C)$ et $\mathcal{GL}_n(\C)$ sont connexes par arcs.
\end{cor}

\begin{thm}
	Le groupe $\mathcal{GL}_n(\R)$ n'est pas connexe par arcs et ses deux composantes sont $\mathcal{GL}_n^+(\R)$ (matrices de déterminant strictement positif) et $\mathcal{GL}_n^-(\R)$ (matrices de déterminant strictement négatif).
\end{thm}

\section{Applications du théorème spectral}

\begin{thm}[Décomposition polaire]
	Toute matrice $A\in\mathcal{GL}_n(\R)$ s'écrit de manière unique sous la forme $A=OS$ avec $O\in\mathcal{O}_n(\R)$ et $S\in\mathcal{S}_n^{++}(\R)$.
\end{thm}

\begin{prop}
	$\mathcal{O}_n(\R)\times\mathcal{S}_n^{++}(\R)$ et $\mathcal{GL}_n(\R)$ sont homéomorphes via la multiplication matricielle.
\end{prop}

\begin{cor}
	Si $G$ est un sous-groupe compact de $\mathcal{GL}_n(\R)$ qui contient $\mathcal{O}_n(\R)$, alors $G=\mathcal{O}_n(\R)$.
\end{cor}

\begin{thm}
	L'exponentielle matricielle induit un homéomorphisme de $\mathcal{S}_n(\R)$ sur $\mathcal{S}_n^{++}(\R)$.
\end{thm}

\section{Décompositions pour les systèmes linéaires}

\begin{dfn}[Déterminants principaux]
	
\end{dfn}

\begin{thm}[Décomposition LU]
	
\end{thm}

\begin{thm}[Critère de Sylvester]
	
\end{thm}

\begin{thm}[Décomposition de Cholesky]
	
\end{thm}

\begin{thm}[Décomposition OT (ou QR)]
	
\end{thm}

\begin{thm}[Décomposition d'Iwasawa]
	
\end{thm}


\end{document}